% mage-auto-heading: prevention-and-observation
% mage-auto-heading: validators-and-gates
% mage-auto-heading: constraining-document-changes
% mage-auto-heading: tests-and-models-work-together
% mage-auto-heading: one-obligation-several-ways-to-hold-it
% mage-auto-heading: one-measurement-different-enforcement-decisions
% mage-auto-heading: positive-and-negative-constraints
% mage-auto-heading: provenance-carried-admission
% mage-auto-heading: the-freedom-a-control-preserves
% mage-auto-heading: what-counts-as-a-mechanism
% Authoritative LaTeX manuscript; Markdown is a generated compatibility projection.

% <!-- part-title: Alignment -->

% <!-- chapter-title: Governance Mechanisms -->

\chapter{Governance Mechanisms}

% <!-- point: four-roles-separate-prevention-observation-judgment-admission | Constraints restrict possibilities, sensors produce evidence, and gates evaluate whether actions or results may proceed and enforce decisions. Checking may be advisory or used within a gate. | terms: governance-mechanism, constraint, sensor, gate -->

Suppose you are adding a reporting service to a billing system. The service may obtain billing data through the billing API, but it must not connect directly to the billing database. Keeping that boundary lets the billing service control access to its data and change its schema without exposing it to every client. An architectural model can state the permitted relationship. How can the engineering environment help preserve it while agents change the reporting code?

% <!-- index-def: governance-mechanism -->

\begin{quote}
\textbf{Governance mechanism.} A repeatable structure in the engineered environment that constrains agentic work, produces evidence about it, or controls whether it may proceed.
\end{quote}

We can distinguish three categories by the engineering questions they answer, as shown in \autoref{governance-mechanism-questions}.

{\footnotesize
\begin{longtable}[]{@{}
  >{\raggedright\arraybackslash}p{(\linewidth - 4\tabcolsep) * \real{0.3333}}
  >{\raggedright\arraybackslash}p{(\linewidth - 4\tabcolsep) * \real{0.3333}}
  >{\raggedright\arraybackslash}p{(\linewidth - 4\tabcolsep) * \real{0.3333}}@{}}
\caption[Three functions of engineering governance]{\emph{Three functions of engineering governance.} Constraints restrict what can happen; sensors provide evidence about what has happened or is proposed; gates determine what may proceed. A single mechanism may perform several functions. This classification draws on established principles of protection and policy enforcement.\citep{saltzer1975protection, rfc2753, oasis2017xacml}}\label{governance-mechanism-questions}\tabularnewline
\toprule\noalign{}
\begin{minipage}[b]{\linewidth}\raggedright
Category
\end{minipage} & \begin{minipage}[b]{\linewidth}\raggedright
Engineering question
\end{minipage} & \begin{minipage}[b]{\linewidth}\raggedright
Effect
\end{minipage} \\
\midrule\noalign{}
\endfirsthead
\toprule\noalign{}
\begin{minipage}[b]{\linewidth}\raggedright
Category
\end{minipage} & \begin{minipage}[b]{\linewidth}\raggedright
Engineering question
\end{minipage} & \begin{minipage}[b]{\linewidth}\raggedright
Effect
\end{minipage} \\
\midrule\noalign{}
\endhead
\bottomrule\noalign{}
\endlastfoot
\textbf{Constraint} & What can happen? & Restricts available actions or states. \\
\textbf{Sensor} & What evidence do we have? & Produces evidence about artifacts, proposed actions, or executions. \\
\textbf{Gate} & May this proceed? & Evaluates whether an action or result may proceed and enforces the decision. \\
\end{longtable}
}

% <!-- box-family: inset -->

% <!-- inset-domain: SOFTWARE ENGINEERING -->

% mage-box-prefix: Inset —

\begin{bookbox}{inset}{Familiar Engineering Machinery}
\label{inset-familiar-engineering-machinery}
\textbf{Constraints} narrow what an implementation or actor can do. Type systems and module boundaries can exclude invalid expressions or hide operations from callers; database constraints and permissions can restrict stored states or resource access. Their assurance depends on the scope they actually enforce. A type-level restriction governs checked code, for example, while inputs from outside that code still need to be checked at the boundary.

\textbf{Sensors} make evidence available for inspection, so their usefulness depends on observability. Logging and tracing expose execution behavior; parsers and static analyses can expose facts about artifacts, configurations, and dependency graphs. The evidence must preserve the distinctions needed by the check, whether it comes from a running system or a static artifact. Collecting it does not itself determine whether work may proceed.

\textbf{Gates} evaluate conditions and enforce decisions at boundaries. A continuous integration (CI) job may report test results; a merge gate uses required results to admit or refuse a change. Deployment gates can require a successful security check, and runtime authorization gates can permit or refuse requests under an access policy. One tool can perform several functions by collecting evidence, evaluating it, and enforcing the resulting decision.
\end{bookbox}

Let\textquotesingle s follow a proposed reporting change from implementation to merge. \autoref{constraint-vs-sensor} shows how module restrictions, test instrumentation, and a merge gate work together to preserve the billing boundary.

\begin{figure}[htbp]
\centering
\includegraphics[width=\linewidth]{book/latex/assets/governance-proposed-change-16bfb894e7.pdf}
\caption{\emph{Governance around a proposed change.} Module restrictions limit the implementation; instrumented tests produce evidence; the merge gate checks the required results before accepting the change. The checks\textquotesingle{} coverage limits the assurance they provide.}
\label{constraint-vs-sensor}
\end{figure}

\section{Prevention and Observation}\label{heading:prevention-and-observation}

% <!-- section-terms: constraints-and-sensors, constraint, sensor -->

% <!-- point: a-mechanism-is-one-of-two-moves-prevent-or-observe | Constraints exclude particular actions or states; sensors supply evidence about behavior that remains possible. | terms: constraint, sensor, constraints-and-sensors -->

% <!-- index-def: constraints-and-sensors -->

% <!-- point: constraint-and-sensor-echo-control-engineering | The accompanying note connects constraints and feedback to systems-safety engineering without identifying agent engineering with a control system. | terms: constraint, sensor, constraints-and-sensors -->

A constraint can exclude some violations before they occur. Sensors provide evidence about behavior that remains possible, so the environment can check it before accepting the change.\mageaside{For the role of constraints and feedback in systems-safety engineering, see Leveson.\citep{leveson2011}}

% <!-- point: a-soft-instruction-becomes-alignment-only-when-the-environment-can-act-on-it | An instruction is enforced only when a mechanism in the environment acts on it. | terms: constraint -->

An instruction to use the billing API can guide an agent, but it does not exclude a direct database connection. Enforcement requires something that changes what the program can express, what it may access, or whether it is accepted. We can start by restricting the available operations.

% <!-- point: a-constraint-prevents-drift-within-the-iteration-and-costs-no-iteration | A constraint excludes particular actions or representations within the scope of its language, interface, or runtime permissions; a closed type does not govern every possible behavior. | terms: constraint, typed-enum -->

% <!-- index-def: constraint -->

% <!-- gloss-only: Constraint | A governance mechanism that restricts available actions or representable states so that particular invalid choices are unavailable through it. A closed enum excludes undeclared state names; transition rules require their own constraint. -->

\textbf{A constraint limits the actions or states the system can express.}\mageaside{Types provide one familiar way to exclude invalid choices. Pierce describes type systems as syntactic methods for ruling out classes of erroneous behavior.\citep{pierce2002} A closed enum with variants such as \texttt{IDLE}, \texttt{RUNNING}, and \texttt{DONE} excludes undeclared state names in checked code; external input must be validated when converted to the type. It does not by itself forbid a transition from \texttt{IDLE} directly to \texttt{DONE}. That transition requires its own constraint.}

For the reporting change, a module boundary exposes the billing API client while keeping the database connector private. The language\textquotesingle s visibility rules prevent checked reporting code from importing that connector. The agent can choose how to assemble the report, but cannot use the private connector to obtain its data. This restriction covers that particular route: code using a generic network library might still attempt a direct database connection.

% <!-- point: selected-structure-can-be-generated-from-a-model | A generator can embody selected obligations in implementation structure; this requires a valid model, a property-preserving transformation, and protection against bypass. | terms: constraint, governance-mechanism -->

Constraints can restrict an implementation after it is produced, or shape its construction so that selected violations cannot arise. Classical model-driven engineering offers a familiar example: a generator translates a model into code, configuration, or other artifacts. If an architectural model determines which services may communicate, generated interfaces or access-control rules can embody those relationships directly. For the properties preserved by the transformation, Alignment is achieved by construction rather than by checking an independently produced implementation afterward.\citep{schmidt2006model,sendall2003model,czarnecki2006feature}

A generator can produce selected parts of an implementation, leaving other decisions to agents. An agent can also customize or extend generated code, provided those changes preserve the obligations embodied in it. The model must express the intended obligation, the generator must preserve it, and subsequent changes must not bypass the generated structure.

% <!-- point: agents-can-construct-checking-and-enforcement-machinery | Agents can construct checkers or controls for explicit obligations; generating a checker does not establish the checked implementation's correctness. | terms: constraint, validator, governance-mechanism -->

Agents can also help construct the checking and enforcement machinery. Once an invariant is explicit, an agent can help translate it into a lint rule, schema constraint, architectural test, or static analysis. Here the generated artifact checks or constrains another implementation. Bodden's work on self-adaptive static analysis proposes automatic optimization of an analysis from an explicit representation of it.\citep{bodden2018,bodden2023sosa} The engineer remains responsible for the obligation and for validating what the resulting mechanism actually establishes.

% <!-- point: a-sensor-observes-and-produces-evidence-and-costs-iterations | A sensor produces evidence about artifacts, attempted actions, or observed executions; observability exposes the state or behavior that the environment must inspect. | terms: sensor, observability -->

% <!-- index-def: sensor -->

% <!-- gloss-only: Sensor | A governance mechanism that observes an artifact, an attempted action, or an executing system and produces evidence. A source scan can record dependencies; runtime instrumentation can record connection attempts and outcomes. That evidence may support checking or other analysis without itself permitting or refusing work. -->

% <!-- index-def: observability -->

\textbf{A sensor observes the work or system and produces evidence.} Sensors depend on \textbf{observability}: the relevant state or behavior must be exposed in a form that the environment can inspect. Logs, traces, instrumentation, and static analyses can supply this evidence.

% <!-- index-example: observability -->

We can run the proposed change in a test environment and record the connections it attempts while processing reporting requests. The instrumentation records the caller, destination, and outcome. When a test exercises a direct database path through a generic network library, the record exposes a violation that the module restriction did not exclude. Recording that connection does not prevent it; the evidence will inform the decision about whether to accept the change.

% <!-- point: a-sensor-is-only-as-good-as-its-observability | Evidence must preserve the distinctions needed to judge the obligation, such as caller identity, destination, and architectural component. | terms: sensor, observability -->

Each connection record must identify the reporting service and distinguish an API call from a direct database connection. A record that says only ``request succeeded'' cannot answer that question. The evidence must preserve the distinctions named by the architectural model before a checker can use it.

\section{Gates: Evaluation and Enforcement}\label{heading:validators-and-gates}

% <!-- section-terms: gate, validator, governance-mechanism -->

% <!-- point: a-gate-supplies-the-consequence-and-a-gate-is-a-commitment | A gate evaluates whether an action or result may proceed, locally or through a delegated check, and enforces the decision; a blocking check can reject both genuine violations and false positives. | terms: gate, validator -->

% <!-- index-def: gate -->

A gate evaluates whether an action or result may proceed and enforces that decision at a boundary. The evaluation may be performed locally or delegated to another checker.

% <!-- point: a-ci-suite-bundles-three-roles-without-collapsing-the-distinction | Instrumented tests produce observations, a checker compares them with the architectural model, and the merge gate requires the build and interaction checks to pass. | terms: sensor, gate, validator, governance-package -->

% <!-- point: a-validator-turns-evidence-into-a-judgment | The merge gate uses build and interaction check results to decide whether a candidate meets its admission conditions. | terms: validator, gate -->

The merge gate requires the build checks to pass and the recorded test interactions to comply with the architectural model. A billing API call is permitted; a direct database connection fails the interaction check, even if the code passed the language\textquotesingle s visibility checks. If a required check fails or its result is missing, the gate refuses the merge and returns the change for repair. Passing these checks establishes that the candidate met those admission conditions; the tests still cover only the executions they exercised.

% <!-- index-def: validator -->

% <!-- gloss-only: Validator | A conventional name for a checking component, such as a schema validator or dependency checker. Its findings may inform a person, produce an advisory report, or supply the evaluation used within a gate. A result does not by itself control whether work proceeds. -->

This example combines a static restriction with dynamic evidence: the language checks imports without running the program, while instrumentation observes test executions. The categories describe each mechanism\textquotesingle s function, not whether it runs before or during execution. A source scanner can supply sensor evidence, and network permissions can constrain connections at runtime. One tool may also perform several functions, such as recording a connection attempt, checking its destination, and refusing it.

% <!-- point: two-roles-act-on-the-work-two-act-at-the-door | Defense in depth combines protections with different coverage; tests and the gate that enforces their results constitute one protection together. | terms: constraint, sensor, gate, validator -->

Combining protections this way can provide \emph{defense in depth}: multiple safeguards reduce dependence on any one of them.\citep{ross2022trustworthysystems} A direct database connection that evades the module restriction may still appear in an instrumented test, causing the merge gate to reject the change. The test and gate form one protection together: the gate enforces the test result but does not independently detect violations. The engineer must consider what each protection covers and which paths could escape both.

% <!-- point: human-process-instantiates-the-same-roles-logging-senses-reviewer-validates-approver-gates | Human review can combine tacit context, evaluation, and approval; selected obligations in agentic work need explicit rules, evidence, and admission consequences. | terms: sensor, gate, constraint, validator -->

A human reviewer may already know that billing owns its database, recognize a forbidden connection in a test trace, and stop a merge. Much of the context and evaluation can remain tacit. When agents carry out repeated changes, engineers must decide which obligations to make explicit, what evidence the environment will collect, and which findings must prevent the work from proceeding. For the reporting service, that means recording the permitted interaction and configuring the check and merge rule, so each new agent need not reconstruct the reviewer's reasoning. Human judgment remains necessary to decide which obligations warrant this treatment.

% <!-- box-family: inset -->

% <!-- inset-domain: CAVEAT -->

% mage-box-prefix: Inset —

\begin{bookbox}{inset}{Enforcement Is Conditional}
\label{inset-enforcement-conditional}
Each check supports a claim within a particular scope. The module restriction excludes imports of the private connector, while the instrumented tests can expose other routes. A trace with no direct database connection establishes only that none was recorded in the observed executions. Other inputs, unobserved processes, or missing instrumentation may reveal different behavior.

Making a check blocking does not broaden its evidence. A test gate requires the selected tests to pass; it does not establish that every untested execution is correct. The architectural model may also omit an important distinction, or the validator may check the wrong obligation. Within its stated scope and assumptions, a gate can make acceptance depend on a condition that the agent cannot simply disregard.
\end{bookbox}

% <!-- point: cheap-evidence-early-full-re-derivation-late-the-thing-asking-to-pass-does-not-define-passing | A release check can regenerate evidence from the actual candidate and evaluate it independently of the producing agent; reuse requires unchanged relevant inputs. | terms: validator, gate -->

A release check can regenerate evidence from the candidate being shipped, rather than trust the producing agent's report or a result from an earlier revision.\mageaside{Repeating the check costs computation and time. Reusing or caching an earlier result can be sound when the system establishes that the relevant inputs have not changed.} This keeps the evidence tied to the artifact being admitted and lets the release check evaluate it independently of the agent that produced the change. The agent's claim that its work passed is not a substitute for that evaluation.

\subsection{Constraining Document Changes}\label{heading:constraining-document-changes}

% <!-- point: typed-document-edits-separate-proposal-from-controlled-application | In DocAble's PDF changeset interface, a closed operation vocabulary constrains proposals, the applier refuses unsupported targets or missing evidence, and subsequent validation checks selected properties of the resulting document. -->

DocAble's PDF changeset interface provides an action boundary for document repair. A reasoning component may decide that an image needs alternative text without receiving unrestricted access to the PDF's internal representation. It proposes a typed operation naming the target and intended change. The operation vocabulary constrains which mutations can be expressed through this interface; the applier determines whether a supported operation can proceed against the current document.

The applier's refusals make that boundary concrete. An operation that marks a structure as decorative must identify the source of the classification. The applier refuses it if the classification source is missing, or if the target is a heading, a formula, or content carrying authored alternative text. These are checks with enforced outcomes. Subsequent validation supplies evidence about selected properties of the resulting document, including content preservation and mechanically testable accessibility requirements. A well-formed proposal does not establish that its description is educationally appropriate, and successful application does not establish complete accessibility.

The boundary also preserves freedom. A specialized pipeline and a general-purpose agent may use different methods to decide what to repair. Neither needs to reproduce the mechanics of applying a supported edit. Nor does either gain authority to bypass the restrictions imposed by the common interface. The architecture is therefore more than a convenient API: it separates the judgment that proposes a change from the mechanisms that determine which changes can be made and which results are acceptable.

\section{Tests and Models Work Together}\label{heading:tests-and-models-work-together}

The reporting-service boundary concerns a relationship between components. Other obligations depend on how state changes over a sequence of operations. Tests already preserve engineering intent about those sequences. An example test records a case and the result it should produce; a regression test keeps checking behavior that a later change might break. Suppose the reporting service queues jobs that workers claim under expiring leases. Worker A's claim expires, worker B acquires a newer claim, and A then tries to release its old one. A regression test can reproduce that sequence against the implementation and assert that B remains the owner. It preserves a concrete failure scenario and makes a recurrence easier to diagnose. Such tests remain useful as the system acquires more explicit models.\citep{winters2020}

We can also state a property covering many such cases: a release carrying an old generation must not clear the current claim. \emph{Property-based testing} supplies generated inputs to a check of that property, varying worker identities, claim generations, and relevant operations. The engineer must choose the property and ensure that the generator exercises the cases that matter, including stale releases. QuickCheck provides a framework for expressing and testing such executable properties.\citep{claessen2000quickcheck}

An explicit model adds a reusable account of the behavior those tests should expect. For this ownership question, the model might retain each job's current owner, claim generation, and whether the claim has expired, together with the effects of claiming, reclaiming, and releasing. A model-based test can use that account to construct sequences of operations and compare the implementation's responses with the permitted outcomes. The model can help choose what to exercise, provide the test oracle---the rule for evaluating the result---or do both. This does not require a second implementation of the database. It requires enough behavior to answer the ownership question.\citep{utting2012taxonomy}

A test driver can issue each operation to both the ownership model and the real lease implementation, then compare their ownership states. An attempted stale release must leave the current claim unchanged in both. The driver translates model operations into real calls and observes their results; the same model supplies expectations across many sequences. A sequential driver will not expose every race inside those calls: concurrent schedules and database atomicity need their own tests or analysis.\citep{prenninger2005abstractions}

Hypothesis's stateful testing can generate sequences of actions and compare a real system with a simpler model, combining property-based and model-based testing. A useful failing sequence can then become a fixed regression test.\citep{hypothesis-stateful}

Model coverage, code coverage, and the properties checked answer different questions. Exercising every modeled transition may miss a problematic ordering; executing every relevant line may still leave ownership unasserted. Passing tests supports a claim only about the executions and properties checked.\citep{winters2020}

The extra model also has a cost. Engineers must validate its rules against the requirements, maintain the connection to the implementation, and revise it when the intended behavior changes. Copying implementation logic into the expected-result calculation can reproduce the same mistake on both sides. A separate model earns its keep when a simpler statement of the behavior makes tests easier to select, results easier to evaluate, or recurring obligations easier to maintain. For a small local rule, a direct assertion may already be sufficient.

\section{One Obligation, Several Ways to Hold It}\label{heading:one-obligation-several-ways-to-hold-it}

% <!-- section-terms: constraint, validator, gate -->

% <!-- point: a-modeled-obligation-does-not-imply-one-enforcement | Generation, runtime constraints, tests, and model checking supply different kinds of restriction or evidence and can be combined for the obligation at hand. -->

A modeled obligation does not imply one particular enforcement technology. Some obligations can be embodied in generated implementation structure; others can be enforced through constraints on agent actions or evaluated through tests and analysis. Tests exercise selected implementation behaviors; model checking examines behavior admitted by an explicit model. Ask what evidence the obligation needs and how the environment should act on it.

% <!-- point: prefer-the-constraint-when-you-can-exclude-the-state | An atomic conditional update can prevent competing claimants from both acquiring the same claim; tests and analysis can examine the surrounding protocol. -->

If an invalid state can be excluded by the runtime operation itself, use that constraint where practical. An atomic conditional update can acquire a job only while its database row remains unclaimed. Two competing workers cannot both win that transition. Tests can still check how the surrounding code handles success, refusal, and failure. The runtime primitive and the tests contribute different protections.

% <!-- point: interleaving-properties-need-a-validator-over-the-state-space | Model checking explores the behavior admitted by an explicit model; model-based testing exercises an implementation against expected model behavior. -->

\emph{Model checking} and \emph{model-based testing} answer different questions. A model checker asks whether the represented protocol permits a violation within the explored domain. Model-based tests ask whether the implementation behaves as the model requires in the executions they perform. Exploring the model does not itself execute database calls or establish that production code follows the modeled transitions. A counterexample found in the model can guide an implementation test, while a test failure can reveal a distinction missing from the model.\citep{baier-katoen2008,utting2012taxonomy}

% <!-- point: temporal-obligations-require-another-shape-of-evidence | A universal eventuality claim requires reasoning about modeled executions and stated scheduling or recovery assumptions; a finite passing test does not establish it. -->

The obligation determines what further evidence is needed. A test can show that a particular job completed, or that it exceeded a required deadline. The stronger claim that every accepted job eventually terminates concerns all executions the system permits. A temporal model checker can examine that claim over a model, subject to its bounds and assumptions about scheduling and recovery. A fairness assumption must be stated rather than silently treating an operation that can occur as one that eventually will. Confidence in the implementation still depends on how faithfully it follows the modeled behavior.\citep{baier-katoen2008,lamport2002}

% <!-- point: recovery-is-not-permission | In DocAble, a lease makes abandoned work recoverable, while a separate durable claim restricts who may perform the protected effect; recovery alone does not supply that permission. | terms: constraint, sensor, constraints-and-sensors -->

% <!-- box-family: inset -->

% <!-- inset-domain: Case study -->

% mage-box-prefix:

\begin{bookbox}{inset}{Worked inset --- DocAble: Recovering a Job Without Repeating Its Effect}
\label{inset-docable-ownership-controls}
DocAble's document-remediation jobs can outlive a worker attempt. A worker may crash, become slow, or lose contact while another attempt becomes eligible to continue the work. A lease makes current ownership visible and allows apparently abandoned work to be recovered. Its expiry cannot establish that the old worker stopped: a delayed worker may still return. The design therefore distinguishes eligibility to recover the work from permission to perform its protected side effect.

As \autoref{match-mechanism-to-property} shows, the runtime claim, the interleaving search, and temporal checking respond to different properties. None substitutes for the others.

\begin{figure}[htbp]
\centering
\includegraphics[width=\linewidth]{book/latex/assets/match-mechanism-to-property-2c4d8f0003.pdf}
\caption{\emph{Different obligations in one job protocol.} An atomic claim excludes competing winners for one transition. Bounded explicit-state exploration checks modeled interleavings. Temporal analysis addresses eventuality under stated assumptions. Tests provide separate evidence about the implementation, and a gate must be chosen separately when an analysis verdict should control whether work may proceed.}
\label{match-mechanism-to-property}
\end{figure}

A separate durable compare-and-set supplies the claim for that effect. It updates the claim only while the stored state matches the required prior state, so competing claimants cannot both acquire the same transition. A repeated delivery or expired lease does not bypass this claim. This is a runtime constraint on a particular effect, rather than a general guarantee that every part of a job executes exactly once. Recovery and permission serve different obligations.

Ownership also depends on the surrounding sequence. After a newer lease generation has been issued, a late release from the old holder must not clear the current claim. DocAble uses a small, purpose-built checker to enumerate reachable interleavings of its ownership model within explicit bounds and search for invariant violations.\mageaside{This is bounded explicit-state model checking. ``Bounded model checking'' conventionally refers to bounded executions encoded for SAT/SMT solving; this checker explores an explicitly bounded state space.} The result concerns the represented protocol. Tests or other analysis must separately establish that the implementation follows it.

A different analysis addresses eventual completion. The system's state and recovery models already carried the requirement that every submitted job eventually reach a terminal state. A small TLA+ specification checked with TLC was then developed for that obligation.\citep{lamport2002} The order matters: engineering identified the required behavior before constructing an analysis of it.

\end{bookbox}

% <!-- point: mechanisms-match-properties-not-an-adoption-sequence | Select and combine mechanisms according to the evidence or restriction they contribute and the cost of obtaining it. -->

Choose each mechanism for the evidence or restriction it contributes and the cost of obtaining it.

% <!-- point: whichever-validator-supplies-the-judgment-the-consequence-depends-on-the-environment | A check may remain advisory; a gate must act on the result when it is intended to control whether work may proceed. -->

A check may inform an engineer without controlling the work. When its result is meant to control whether work may proceed, a gate must use it to permit or refuse the action or artifact.

\section{One Measurement, Different Enforcement Decisions}\label{heading:one-measurement-different-enforcement-decisions}

% <!-- section-terms: constraint, sensor, gate -->

% <!-- point: one-model-several-enforcement-decisions | One quantitative model can support observation, adaptation, graceful degradation, or hard admission control; the appropriate response depends on the obligation attached to the quantity. | terms: constraint, sensor, gate -->

The reporting service's usage measurements can support several responses, depending on the policy attached to them. An engineer may first record API cost for accounting. The service may reduce optional report detail as a job approaches its budget, or refuse new reports when a configured daily ceiling would be exceeded. These are different decisions about what the environment should permit, even when they draw on related measurements. \autoref{cost-enforcement-gradient} shows observation, adaptation, and admission control as alternative responses. The appropriate response follows from the obligation, rather than from the mere availability of a measurement.

\begin{figure}[htbp]
\centering
\includegraphics[width=\linewidth]{book/latex/assets/cost-enforcement-gradient-4e2626b031.pdf}
\caption{\emph{One Measurement, Different Enforcement Decisions.} The same measured quantity may be observed without enforcement, used to adapt behavior, trigger graceful degradation, or control admission. The appropriate response follows from the obligation attached to the measurement.}
\label{cost-enforcement-gradient}
\end{figure}

% <!-- point: a-quantitative-tolerance-can-expose-margin | A quantitative tolerance can expose margin — the room between the realized value and its bound; two conforming realizations can differ in how much room they keep. | terms: tolerance, sensor -->

A quantitative tolerance can also expose margin. If a value remains inside its acceptable bound, margin describes the room between the realized value and that boundary. A reporting request completing in 120 ms against a 200 ms ceiling has 80 ms of margin. That information can matter before conformance fails: two realizations may both satisfy the same obligation while one operates much closer to its boundary and therefore has less room to absorb variation or future change.

% <!-- point: margin-needs-a-meaningful-distance | Margin is computable only where the model supplies a meaningful distance; many software obligations are qualitative and admit none. | terms: tolerance -->

That calculation is easy only when the model supplies a meaningful distance to the boundary. Quantitative obligations often do: latency, memory use, throughput, or cost can be compared directly with a declared bound. Many software obligations are qualitative instead. A dependency may cross a forbidden architectural boundary or remain within the permitted graph; a transition may be legal or illegal. Some such obligations admit useful notions of distance from failure, while others do not. Margin is useful where the model provides a meaningful measure; qualitative obligations may need a different way to assess risk.

\section{Positive and Negative Constraints}\label{heading:positive-and-negative-constraints}

% <!-- section-terms: constraint -->

% <!-- point: a-constraint-can-forbid-an-action-or-require-an-evidence-producing-one | A constraint can forbid an action or require an evidence-producing one; provenance shows the positive form. | terms: constraint -->

The reporting-service rule forbids a particular access. Constraints can also require an accompanying action, such as recording who performed a mutation. A document-editing interface, for example, might prohibit callers from using the raw format library directly. It might also require each permitted mutation to emit an attribution record. A source-time rule can check that the recording call appears in each known mutation path. \autoref{positive-negative-obligation} sets these two obligations side by side: restrict the prohibited access, and require the evidence-producing behavior. Checking that the call is present still leaves the correctness of the emitted evidence to its own checks.

\begin{figure}[htbp]
\centering
\includegraphics[width=\linewidth]{book/latex/assets/positive-negative-obligation-59f450f724.pdf}
\caption{\emph{Negative and positive constraints.} A negative constraint removes a forbidden action; a positive constraint requires a specified action or evidence-producing behavior. Both restrict the admissible implementation space.}
\label{positive-negative-obligation}
\end{figure}

\section{Provenance-Carried Admission}\label{heading:provenance-carried-admission}

% <!-- section-terms: constraint, governance-mechanism -->

% <!-- point: provenance-carried-admission-is-the-derivation-path-case-of-conformance | Provenance is useful for a particular architectural-conformance problem: the required relation is the path through which a protected operation was derived, not merely which components interact. | terms: constraint -->

The reporting-service example concerns which components may interact. Some architectural obligations also concern the path through which an operation was derived. They require evidence that a necessary step occurred before the protected effect.

% <!-- point: restrict-a-recognizable-primitive-surface-directly | Where a consequential effect passes through a recognizable primitive, restrict that surface directly — the prohibited interaction has a tractable signature. | terms: constraint -->

The reporting boundary can be enforced by controlling a recognizable primitive: the connection to the billing database. More generally, when the environment can identify the APIs that perform a consequential effect and designate the components permitted to use them, it can check known source-level callers and restrict access at runtime. The mechanism must cover the routes through which the protected effect can actually occur.

% <!-- point: some-obligations-have-no-single-forbidden-call | Some obligations have no single forbidden call; a protected operation is reachable by composing otherwise legitimate APIs into a new path. | terms: constraint -->

Other architectural obligations do not have such a convenient boundary. Suppose a payment must be approved before it reaches the transfer operation. A future implementation may combine otherwise legitimate account queries, calculations, and service calls into a new route to that operation. Every individual call can be permitted while the route omits approval. A rule looking only for previously identified bypasses may therefore miss the new path.

An agent may find an implementation path its designers did not anticipate. If the path bypasses approval, a rule written to catch an earlier bypass may miss it. Trying to enumerate every forbidden route can therefore leave gaps as the implementation changes.

% <!-- point: provenance-reverses-the-enforcement-problem | Provenance reverses enforcement: require the protected operation to carry evidence that it passed through the sanctioned abstraction. | terms: constraint, governance-mechanism -->

Provenance can reverse the enforcement problem. Instead of trying to recognize every way an operation might have bypassed the sanctioned abstraction, require the protected operation to carry evidence that it passed through that abstraction.

In the payment example, the approval service can issue a digitally signed token that cryptographically binds its approval to the proposed transfer: its source account, destination, and amount. Callers without the signing key cannot forge an approval or change those fields while retaining a valid signature. The transfer boundary verifies the signature and checks that the approved details match the requested operation before admitting it. The boundary must also atomically consume a single-use approval so concurrent requests cannot reuse it for a duplicate transfer.\citep{owasp-transaction-authorization} The general admission condition is:

\[
\operatorname{admit}(x) \iff \operatorname{valid}(x) \land \operatorname{sanctionedProvenance}(x)
\]

\autoref{provenance-carried-admission} traces the single sanctioned path: an operation reaches the protected boundary only by deriving through the sanctioned abstraction, which supplies the provenance the boundary demands.

\begin{figure}[htbp]
\centering
\includegraphics[width=\linewidth]{book/latex/assets/provenance-carried-admission-835241521f.pdf}
\caption{\emph{Provenance-carried admission.} A protected operation must derive through the sanctioned abstraction, which attaches provenance to the resulting operation. The receiving boundary requires that provenance before admitting the operation. The mechanism removes the degree of freedom to obtain the operation through some other implementation path, without requiring the environment to enumerate those alternative paths.}
\label{provenance-carried-admission}
\end{figure}

An agent remains free to invent new internal implementations, but invention does not create a new route around the obligation. A newly synthesized path that bypasses the sanctioned abstraction arrives at the boundary without valid provenance and is refused. Enforcement therefore depends on recognizing the allowed derivation, not enumerating all possible forbidden derivations.

% <!-- point: choose-provenance-when-derivation-paths-are-numerous-or-compositional | Restrict a small recognizable surface directly; reach for provenance when the obligation concerns how an ordinary operation was derived and its paths are numerous or compositional. | terms: constraint -->

If the protected effect passes through a small, recognizable set of operations, restricting those operations may be sufficient. Provenance becomes useful when the obligation concerns how an operation was prepared and the alternative routes are too numerous or difficult to check directly. The engineer can then require evidence of the approved preparation at the point where the effect occurs.

% <!-- point: provenance-makes-a-derivation-path-sanctioned | Provenance here is not attribution; it makes a derivation path the sanctioned route to admission — producer-side rules check known emitters, the consumer boundary refuses unknown routes. | terms: constraint, governance-mechanism -->

In the payment example, provenance establishes that the required approval occurred before the transfer. A source-time rule can check that known implementations attach the approval token. The receiving boundary can require a valid token for every transfer, including requests from implementations the source rule did not inspect. One check examines known producers; the other protects the receiving operation.

The control restricts a derivation path through the interaction graph: the protected operation must pass through the node that issues its provenance. Other implementation choices remain open.

\section{The Freedom a Control Preserves}\label{heading:the-freedom-a-control-preserves}

% <!-- section-terms: constraint, governance-mechanism -->

% <!-- point: preserving-freedom-has-positive-value | Preserving degrees of freedom has positive value: some choices are immaterial, some are not yet ripe, and agent exploration can find strategies engineers did not anticipate. | terms: constraint, governance-mechanism -->

An engineer need not fix every choice that a control could restrict. Some choices are immaterial to the obligation: several implementations satisfy it equally, and fixing one would only add coupling. Other choices are consequential but not yet understood well enough to settle. Leaving them open gives implementation and exploration a chance to supply more evidence; an agent may find a strategy its engineers did not anticipate.

As implementation becomes cheaper, such exploration becomes increasingly attractive. An engineer can specify what must hold, allow an agent to investigate how to achieve it, and commit to a mechanism once the evidence justifies the choice. A well-designed control excludes unacceptable realizations without unnecessarily restricting the search among acceptable ones.

% <!-- point: defer-consequential-commitments-until-the-last-responsible-moment | Defer consequential commitments until the last responsible moment; cheap implementation changes the economics of waiting, and enforcement makes the deferral safer. | terms: constraint -->

This is the familiar discipline of deferring a consequential commitment until the \emph{last responsible moment}. Waiting can yield useful information, but it can also make a later change more expensive. Cheaper agent implementation can shorten that wait by producing a prototype or experiment sooner. Controls that protect known obligations can make such exploration safer, allowing engineers to leave more of the implementation open while they investigate.

% <!-- point: realization-space-underconstraint-vs-overconstraint | Obligations accept the region A, controls permit the region P; underconstraint is P minus A, overconstraint is A minus P. | terms: constraint, governance-mechanism -->

We can describe this tradeoff with sets. Let \(\mathcal R\) be the possible realizations of the work. Each obligation \(O_i \subseteq \mathcal R\) is the set of realizations satisfying it, so the acceptable realizations are \(\mathcal A = O_1 \cap \ldots \cap O_n\).

For a fixed control configuration and decision context, let \(C_j \subseteq \mathcal R\) be the set of realizations permitted by enforcing control \(j\), whether a direct constraint or a gate. When every applicable control must permit a realization, the environment permits \(\mathcal P = C_1 \cap \ldots \cap C_m\). A gate can exclude a realization that satisfies all direct constraints, for example because required evidence is missing. Sensors can supply evidence for that decision without themselves restricting \(\mathcal P\). The permitted region can change as the evidence or controls change. Two kinds of mismatch appear as set differences:

\[
\text{underconstraint: } \mathcal P \setminus \mathcal A \qquad \text{overconstraint: } \mathcal A \setminus \mathcal P
\]

Underconstraint leaves unacceptable realizations permitted; it threatens correctness. Overconstraint excludes acceptable realizations; it sacrifices useful engineering freedom. \autoref{permitted-vs-acceptable} draws both regions and their mismatch.

\begin{figure}[htbp]
\centering
\includegraphics[width=\linewidth]{book/latex/assets/permitted-vs-acceptable-fc09428d74.pdf}
\caption{\emph{Obligations and permitted realizations.} The obligations jointly admit an acceptable region; the enforcing controls jointly permit another. Where permitted exceeds acceptable, unacceptable realizations slip through --- underconstraint; where acceptable exceeds permitted, acceptable alternatives are excluded --- overconstraint.}
\label{permitted-vs-acceptable}
\end{figure}

% <!-- point: matching-permitted-to-acceptable-is-an-idealization-restriction-must-be-justified | Making the permitted region coincide with the acceptable one is an idealization; a deliberate restriction is legitimate — the point is a justified decision, not an accidental consequence. | terms: constraint -->

Making \(\mathcal P\) coincide with \(\mathcal A\) is an idealization, not always a practical or even desirable goal. An engineer may deliberately exclude acceptable alternatives --- for operational simplicity, for consistency, or for assurance. Comparing the two regions helps engineers identify that restriction and decide whether its benefit justifies the lost freedom.

\section{What Counts as a Mechanism}\label{heading:what-counts-as-a-mechanism}

% <!-- section-terms: governance-mechanism, constraints-and-sensors -->

% <!-- point: keep-policy-guidance-and-governance-mechanism-separate | Keep policy, guidance, and governance mechanism separate; enforcement begins when something acts. | terms: governance-mechanism, model-as-map -->

\textbf{Policy}, \textbf{guidance}, and \textbf{governance mechanisms} contribute different things to engineering work. A model can state a policy, and a playbook can explain how to follow it. Sensor evidence and check results can help engineers assess the work without changing what the environment permits. Enforcement requires a constraint on the relevant action or a gate that evaluates a condition and permits or refuses passage.

% <!-- point: the-same-artifact-can-play-several-roles | One modeled policy can guide reasoning, support checking and admission, and generate runtime permissions. | terms: governance-mechanism, model-as-map -->

For the reporting service, an agent can consult the architectural model when choosing how to obtain billing data. The dependency validator uses the same model to check the relationships extracted from the code. A build gate can require that check to pass. A deployment tool can also derive network permissions from the model, allowing API calls while denying direct database connections. Each use needs its own implementation: supplying the model to the agent does not install the validator or the permissions.

% <!-- point: implementations-bundle-roles-name-them-do-not-stretch-the-terms | A governance package can combine models, sensors, constraints, and gates with the checks used to evaluate evidence. | terms: governance-package, constraints-and-sensors -->

% <!-- gloss-only: Governance mechanism | A repeatable structure in the engineered environment that constrains agentic work, produces evidence about it, or controls whether it may proceed. Constraints restrict available actions or states; sensors produce evidence; gates evaluate whether actions or results may proceed and enforce the resulting decisions. A control is a governance mechanism considered in relation to the obligation or risk it addresses. -->

% <!-- index-def: governance-package -->

% <!-- gloss-only: Governance package | A maintained combination of models, mechanisms, and supporting checks. For example, a package may include an architectural model, a dependency scanner, a checker that compares the extracted dependencies with the model, and a merge gate that uses the result to admit or refuse a change. -->

A \textbf{governance package} combines these mechanisms so they can be maintained and used together. Within the package, keep track of which mechanism produces evidence, which evaluates it, and which restricts or refuses work. That distinction makes it easier to diagnose a failure without confusing an incorrect judgment with a failure to enforce it.

% <!-- point: represented-intent-becomes-enforceable-through-one-of-four-roles | Constraints restrict possibilities; sensors supply evidence; gates evaluate whether actions or results may proceed and enforce decisions. The next question is where these mechanisms come from. | terms: governance-mechanism, alignment-principle -->

Some mechanisms can be designed before the first autonomous change; others become visible only after the environment fails. The next section shows how those failures become durable structure.
